% GENERATED FILE. Edit canonical lessons, then run npm run generate:books.
% canonical-source-hash: fnv1a64:44599c0001373d5c
% canonical-lessons: LA-C171-fallo, LA-C171-tolero, LA-C171-condo, LA-C171-creo, LA-C171-credo

\chapter{To deceive, To endure, To found, To create, To believe}
\label{ch:la-to-deceive-to-endure-to-found-to-create-to-believe}
\hlchaptermodality{171}

\begin{chapteropening}
\textbf{By the end of this chapter:} \emph{I can say to deceive, to endure, to found, to create and to believe.}

Complete the last lesson of chapter 171: I can say to deceive, to endure, to found, to create and to believe.
\end{chapteropening}

\section[fallō, fallere]{fallō, fallere — to deceive}

\label{lesson:LA-C171-fallo}



\begin{quote}
Before the new one: say the Latin for to suspect, then the Latin for to press.
\end{quote}

\subsection*{You'll want to know: fallō, fallere}
\textbf{fallō, fallere} — \textquotedblleft{}to deceive\textquotedblright{}.

\begin{cousinweb}
\textbf{fallō} has no sure origin; Greek \emph{sphallō}, \textquotedblleft{}to trip up\textquotedblright{} (Lewis \& Short), is often compared, and EtymDB 2.1 records a proposed root \emph{*ghwel-}, \textquotedblleft{}to deceive\textquotedblright{}.
\end{cousinweb}

\begin{grammarlens}[title={What you've built}]
One more word for this chapter.
\end{grammarlens}

\subsection*{Your turn}
\begin{itemize}
\raggedright
  \item \emph{Say it:} \emph{fallō, fallere}
  \item \emph{Say it:} \emph{fallō, fallere}, once more
  \item \emph{Say it:} say \emph{fallō, fallere}
  \item \emph{Recall:} say the Latin for to suspect, then the Latin for to press, then say \emph{fallō, fallere} again
\end{itemize}

\begin{tcolorbox}[breakable,colback=teal!4,colframe=teal!35!black,title={Before you move on}]
What does fallō, fallere mean? (\textquotedblleft{}To deceive\textquotedblright{}.) Say it once more.
\end{tcolorbox}
\section[tolerō, tolerāre]{tolerō, tolerāre — to endure, to put up with}

\label{lesson:LA-C171-tolero}



\begin{quote}
Before the new one: say the Latin for to press, then the Latin for to deceive.
\end{quote}

\subsection*{You'll want to know: tolerō, tolerāre}
\textbf{tolerō, tolerāre} — \textquotedblleft{}to endure, to put up with\textquotedblright{}.

\begin{cousinweb}
\textbf{tolerō} is inherited, usually traced to Proto-Indo-European \emph{*telh2-}, \textquotedblleft{}to bear, to endure\textquotedblright{} (EtymDB 2.1); \emph{tollō}, \textquotedblleft{}to lift\textquotedblright{}, comes from the same root.
\end{cousinweb}

\begin{grammarlens}[title={What you've built}]
One more word for this chapter.
\end{grammarlens}

\subsection*{Your turn}
\begin{itemize}
\raggedright
  \item \emph{Say it:} \emph{tolerō, tolerāre}
  \item \emph{Say it:} \emph{tolerō, tolerāre}, once more
  \item \emph{Say it:} say \emph{tolerō, tolerāre}
  \item \emph{Recall:} say the Latin for to press, then the Latin for to deceive, then say \emph{tolerō, tolerāre} again
\end{itemize}

\begin{tcolorbox}[breakable,colback=teal!4,colframe=teal!35!black,title={Before you move on}]
What does tolerō, tolerāre mean? (\textquotedblleft{}To endure, to put up with\textquotedblright{}.) Say it once more.
\end{tcolorbox}
\section[condō, condere]{condō, condere — to found; to hide away}

\label{lesson:LA-C171-condo}



\begin{quote}
Before the new one: say the Latin for to deceive, then the Latin for to endure.
\end{quote}

\subsection*{You'll want to know: condō, condere}
\textbf{condō, condere} — \textquotedblleft{}to found; to hide away\textquotedblright{}.

\begin{cousinweb}
\textbf{condō} is \emph{con-}, \textquotedblleft{}together\textquotedblright{}, with \emph{-dō}, \textquotedblleft{}to put\textquotedblright{} (Lewis \& Short): to put together, to found, or to put away. This \emph{-dō} is usually traced to Proto-Indo-European \emph{*dheh1-}, \textquotedblleft{}to put, to place\textquotedblright{} (EtymDB 2.1), and merged in form with \emph{dō}, \textquotedblleft{}to give\textquotedblright{}.
\end{cousinweb}

\begin{grammarlens}[title={What you've built}]
One more word for this chapter.
\end{grammarlens}

\subsection*{Your turn}
\begin{itemize}
\raggedright
  \item \emph{Say it:} \emph{condō, condere}
  \item \emph{Say it:} \emph{condō, condere}, once more
  \item \emph{Say it:} say \emph{condō, condere}
  \item \emph{Recall:} say the Latin for to deceive, then the Latin for to endure, then say \emph{condō, condere} again
\end{itemize}

\begin{tcolorbox}[breakable,colback=teal!4,colframe=teal!35!black,title={Before you move on}]
What does condō, condere mean? (\textquotedblleft{}To found; to hide away\textquotedblright{}.) Say it once more.
\end{tcolorbox}
\section[creō, creāre]{creō, creāre — to create, to elect}

\label{lesson:LA-C171-creo}



\begin{quote}
Before the new one: say the Latin for to endure, then the Latin for to found; to hide away.
\end{quote}

\subsection*{You'll want to know: creō, creāre}
\textbf{creō, creāre} — \textquotedblleft{}to create, to elect\textquotedblright{}.

\begin{cousinweb}
\textbf{creō} is inherited from Proto-Italic \emph{*krēāō}, \textquotedblleft{}to make grow\textquotedblright{}, usually traced to Proto-Indo-European \emph{*ker-}, \textquotedblleft{}to grow\textquotedblright{} (EtymDB 2.1); \emph{crēscō}, \textquotedblleft{}to grow\textquotedblright{}, is its partner. English \emph{create} comes from it.
\end{cousinweb}

\begin{grammarlens}[title={What you've built}]
One more word for this chapter.
\end{grammarlens}

\subsection*{Your turn}
\begin{itemize}
\raggedright
  \item \emph{Say it:} \emph{creō, creāre}
  \item \emph{Say it:} \emph{creō, creāre}, once more
  \item \emph{Say it:} say \emph{creō, creāre}
  \item \emph{Recall:} say the Latin for to endure, then the Latin for to found; to hide away, then say \emph{creō, creāre} again
\end{itemize}

\begin{tcolorbox}[breakable,colback=teal!4,colframe=teal!35!black,title={Before you move on}]
What does creō, creāre mean? (\textquotedblleft{}To create, to elect\textquotedblright{}.) Say it once more.
\end{tcolorbox}
\section[crēdō, crēdere]{crēdō, crēdere — to believe, to trust}

\label{lesson:LA-C171-credo}



\begin{quote}
Before the new one: say the Latin for to found; to hide away, then the Latin for to create.
\end{quote}

\subsection*{You'll want to know: crēdō, crēdere}
\textbf{crēdō, crēdere} — \textquotedblleft{}to believe, to trust\textquotedblright{}.

\begin{cousinweb}
\textbf{crēdō} is inherited, usually traced to a Proto-Indo-European phrase \emph{*kred dheh1-}, \textquotedblleft{}to put one's heart\textquotedblright{} (EtymDB 2.1), from \emph{*kerd-}, \textquotedblleft{}heart\textquotedblright{}, and \emph{*dheh1-}, \textquotedblleft{}to put\textquotedblright{}.
\end{cousinweb}

\begin{grammarlens}[title={What you've built}]
One more word for this chapter.
\end{grammarlens}

\subsection*{Your turn}
\begin{itemize}
\raggedright
  \item \emph{Say it:} \emph{crēdō, crēdere}
  \item \emph{Say it:} \emph{crēdō, crēdere}, once more
  \item \emph{Say it:} say \emph{crēdō, crēdere}
  \item \emph{Recall:} say the Latin for to found; to hide away, then the Latin for to create, then say \emph{crēdō, crēdere} again
\end{itemize}

\begin{tcolorbox}[breakable,colback=teal!4,colframe=teal!35!black,title={Before you move on}]
What does crēdō, crēdere mean? (\textquotedblleft{}To believe, to trust\textquotedblright{}.) Say it once more.
\end{tcolorbox}
